\section*{Bloque C. Casos con implicancias regulatorias y laborales}

\noindent\textit{Preguntas 66 a 90. Aplican el análisis de los mercados de factores a las empresas de agua y saneamiento, la demanda laboral, el monopsonio en zonas rurales, los sindicatos, el salario mínimo y la negociación bilateral.}

\q{La demanda de trabajo de una empresa de agua es derivada porque:}
\begin{opts}
  \item depende de la demanda de servicios de agua y de la productividad del trabajo.
  \item la determinan únicamente los sindicatos.
  \item es fija e independiente del servicio.
  \item no depende del salario.
\end{opts}

\q{Si aumenta la demanda de servicios de agua, la demanda de trabajo de la empresa:}
\begin{opts}
  \item disminuye.
  \item aumenta.
  \item no cambia.
  \item se vuelve cero.
\end{opts}

\q{Si la empresa sustituye la lectura manual de medidores por tecnología digital, la demanda de trabajo para esa tarea:}
\begin{opts}
  \item aumenta.
  \item no cambia.
  \item se vuelve cero.
  \item disminuye.
\end{opts}

\q{En una localidad pequeña, la empresa de agua es el único empleador de técnicos. Esta situación configura:}
\begin{opts}
  \item un monopolio.
  \item competencia perfecta.
  \item un monopsonio.
  \item un monopolio bilateral.
\end{opts}

\q{El monopsonista paga un salario:}
\begin{opts}
  \item igual al MRP.
  \item mayor que el MRP.
  \item igual a cero.
  \item menor que el MRP.
\end{opts}

\q{Si se fija un salario mínimo entre el salario de monopsonio y el competitivo, el empleo:}
\begin{opts}
  \item disminuye.
  \item no cambia.
  \item aumenta.
  \item se vuelve cero.
\end{opts}

\q{Si el salario mínimo se fija por encima del salario competitivo, el empleo:}
\begin{opts}
  \item aumenta.
  \item disminuye.
  \item no cambia.
  \item se vuelve cero.
\end{opts}

\q{En una negociación colectiva entre la empresa de agua y el sindicato, el salario resultante depende de:}
\begin{opts}
  \item solo la demanda de trabajo.
  \item solo la oferta de trabajo.
  \item la tecnología.
  \item el poder de negociación de las partes.
\end{opts}

\q{En un monopolio bilateral, el salario queda:}
\begin{opts}
  \item en cero.
  \item por encima del MRP.
  \item igual al competitivo.
  \item dentro de un rango delimitado por el salario de reserva y el salario máximo que la empresa está dispuesta a pagar.
\end{opts}

\q{Si el sindicato prioriza el empleo, preferirá:}
\begin{opts}
  \item un salario menor y mayor empleo.
  \item un salario mayor y menor empleo.
  \item un salario cero.
  \item no negociar.
\end{opts}

\q{Una huelga en una empresa de agua:}
\begin{opts}
  \item impone costos a la empresa y afecta el poder de negociación.
  \item es un subsidio.
  \item es un impuesto al trabajo.
  \item es un precio máximo.
\end{opts}

\q{El arbitraje obligatorio en un conflicto laboral:}
\begin{opts}
  \item aumenta la huelga.
  \item resuelve el conflicto mediante una decisión de un tercero.
  \item elimina el sindicato.
  \item fija el precio del agua.
\end{opts}

\q{La capacitación del personal de la empresa de agua:}
\begin{opts}
  \item aumenta la productividad y puede elevar los salarios.
  \item reduce la productividad.
  \item no tiene efectos.
  \item solo genera costos.
\end{opts}

\q{La inversión en infraestructura que aumenta la productividad del trabajo:}
\begin{opts}
  \item reduce la demanda de trabajo.
  \item no afecta la demanda de trabajo.
  \item aumenta la demanda de trabajo.
  \item la vuelve cero.
\end{opts}

\q{La tercerización de servicios en una empresa de agua:}
\begin{opts}
  \item siempre aumenta el empleo total.
  \item puede cambiar la composición de la demanda de trabajo.
  \item no tiene efecto.
  \item elimina todo el empleo.
\end{opts}

\q{Una reducción de ingresos de la empresa por una regulación tarifaria estricta puede:}
\begin{opts}
  \item reducir la demanda de trabajo o moderar los salarios.
  \item aumentar automáticamente el empleo.
  \item no afectar el empleo.
  \item eliminar la demanda de agua.
\end{opts}

\q{Un diferencial salarial que compensa la lejanía y las condiciones difíciles de una zona rural es:}
\begin{opts}
  \item una discriminación de precios.
  \item un diferencial salarial compensatorio.
  \item un subsidio al empleo.
  \item un salario mínimo.
\end{opts}

\q{La oferta de trabajo para técnicos en zonas remotas suele ser:}
\begin{opts}
  \item perfectamente elástica.
  \item horizontal.
  \item cero.
  \item relativamente inelástica.
\end{opts}

\q{Si la oferta de trabajo es inelástica, un aumento de la demanda de trabajo:}
\begin{opts}
  \item no afecta el salario.
  \item reduce el salario.
  \item eleva el salario.
  \item lo vuelve cero.
\end{opts}

\q{La demanda de trabajo de la empresa de agua es más elástica cuando:}
\begin{opts}
  \item no existen sustitutos del trabajo.
  \item el costo laboral es una proporción importante y hay sustitutos disponibles.
  \item el salario es fijo.
  \item la oferta de trabajo es vertical.
\end{opts}

\q{Si el salario aumenta y la empresa no puede sustituir trabajo, el empleo:}
\begin{opts}
  \item cae menos.
  \item cae más.
  \item aumenta.
  \item se vuelve cero.
\end{opts}

\q{La negociación bilateral puede llevar a un contrato eficiente si las partes negocian:}
\begin{opts}
  \item solo el salario.
  \item sin información.
  \item únicamente el gobierno.
  \item salario y empleo conjuntamente.
\end{opts}

\q{Si el sindicato tiene poder de monopolio y la empresa poder de monopsonio, el resultado salarial es:}
\begin{opts}
  \item el competitivo.
  \item cero.
  \item el de monopolio.
  \item indeterminado y depende de la negociación.
\end{opts}

\q{Una regulación laboral que eleva el costo de despido:}
\begin{opts}
  \item aumenta la contratación.
  \item no tiene efectos.
  \item reduce la flexibilidad para ajustar el empleo.
  \item elimina el mercado de trabajo.
\end{opts}

\q{La evaluación de una política salarial en una empresa de agua debe considerar:}
\begin{opts}
  \item solo el salario.
  \item solo los beneficios.
  \item solo al sindicato.
  \item efectos en el empleo, los costos, la calidad del servicio y la equidad.
\end{opts}
